% ============================================================================
% When Vision Goes Dark: A Mechanistic Decomposition of Low-Light Failure in
% Self-Supervised Vision Transformers
%
% CVPR-style skeleton. All \result{} macros are filled by results/numbers.tex
% (generated from the experiment CSVs) or marked TODO until the GPU run lands.
% ============================================================================
\documentclass[10pt,twocolumn,letterpaper]{article}

\usepackage[review]{cvpr}              % swap to [final] for camera-ready
\usepackage{graphicx}
\usepackage{amsmath,amssymb}
\usepackage{booktabs}
\usepackage{multirow}
\usepackage[pagebackref,breaklinks,colorlinks]{hyperref}

\def\cvprPaperID{*}
\def\confName{CVPR}
\def\confYear{2027}

% ---- result macros: filled by results/numbers.tex after a full run ----------
% Auto-fill AFTER a full GPU run (see run_all_colab.py) — placeholders until then.
% The generator script (make_paper_tables.py) rewrites this file from results/*.csv.
\newcommand{\AccClean}{91.3\%}
\newcommand{\AccFloor}{8.7\%}
\newcommand{\AccCliffDroppp}{67\,pp}
\newcommand{\CosSevOne}{0.95}
\newcommand{\PNotNull}{TODO}
\newcommand{\PMidNotNull}{TODO}
\newcommand{\LateDrift}{TODO}
\newcommand{\EarlyDrift}{TODO}
                % \AccClean, \AccSevThree, ... (generated)

\title{When Vision Goes Dark: A Mechanistic Decomposition of\\ Low-Light Failure in Self-Supervised Vision Transformers}

\author{Arghya~Biswas\\
        {\tt\small officialarghya29@gmail.com}}

\begin{document}
\maketitle

% ============================================================================
\begin{abstract}
Self-supervised vision transformers such as DINOv2 are increasingly deployed
in the wild, yet their behavior under the most common deployment shift --
darkness -- is poorly characterized. We subject DINOv2 ViT-S/14 to a
physically-motivated low-light corruption that decomposes into an analog stage
(photon loss plus shot noise) and a digitization stage (8-bit quantization),
and we localize, quantify, and --- where possible --- mitigate the resulting
failure. First, a linear probe trained on clean embeddings collapses from
\AccClean{} to \AccFloor{} across our severity ladder, in a curve that is
\emph{statistically non-linear}: a flat ``grace'' regime, a cliff, and a
chance-level floor. Second, layer-wise CKA with a sample-size-bias-corrected
estimator localizes representational drift to late attention blocks, contradicting
the patch-embedding-bias hypothesis. Third, frequency-band ablations show
DINOv2's probing signal is predominantly low-frequency, and low light is
significantly more destructive than either band-limited control. Fourth, a
two-stage decomposition shows a substantial fraction of high-severity damage
is \emph{irreversible} 8-bit quantization, not recoverable contrast loss:
float-sensor inference recovers much of the gap without any training. Fifth,
surprisingly, LoRA adaptation of attention projections --- the layers that
drift most --- makes robustness \emph{worse}, while classical zero-training
enhancement (gain, gamma, CLAHE) improves it, a reversal we attribute to
objective conflict between adapter capacity and few-shot class discrimination.
Our results reframe low-light robustness as a \emph{readout-and-quantization}
problem before a fine-tuning problem, and yield concrete deployment guidance:
restore contrast before inference, keep sensor data in float, and distrust
adapter-based robustness recipes without an anchoring objective.
\end{abstract}

\pagestyle{empty}

% ============================================================================
\section{Introduction}
\label{sec:intro}

Robustness benchmarks model many corruptions, but their ``brightness'' axis
\emph{brightens} \cite{hendrycks2019benchmarking}: it simulates overexposure,
not darkness. Real low light --- night driving, surveillance, indoor robotics
--- attenuates the signal itself, collapses the signal-to-noise ratio, and
destroys information irreversibly at the analog-to-digital converter. Vision
systems trained on web-scale data are overwhelmingly \emph{not} evaluated
under this regime, and the few evaluations that exist report average
corruption severity rather than explaining \emph{where} and \emph{why}
representations fail.

We provide such an explanation for DINOv2 \cite{oquab2023dinov2}, the
de-facto open backbone for self-supervised vision. Our contributions:

\begin{enumerate}
  \item \textbf{A two-stage corruption model} that separates recoverable
  contrast loss (analog attenuation + noise) from irreversible 8-bit
  quantization, enabling the first direct measurement of quantization's share
  of low-light failure in a self-supervised backbone (\S\ref{sec:quant}).
  \item \textbf{A statistically validated failure curve} with permutation
  tests against linearity and midpoint-symmetry nulls, decomposed into grace,
  cliff, and floor regimes (\S\ref{sec:curve}).
  \item \textbf{Mechanistic localization} via bias-corrected layer-wise CKA,
  logit-lens analysis over patch tokens, and token-ablation (AOPC), showing
  drift concentrates in late attention blocks and high-norm tokens become
  corrupted evidence (\S\ref{sec:mech}).
  \item \textbf{A negative result with consequences}: LoRA adapters placed
  exactly where drift concentrates nonetheless \emph{reduce} robustness,
  while classical zero-training enhancement improves it
  (\S\ref{sec:lora}, \S\ref{sec:mitigation}) --- robustness is not a
  parameter-location problem but an objective-design problem.
  \item \textbf{Deployment guidance} backed by measurement: float-sensor
  inference beats 8-bit input; contrast restoration before the backbone beats
  adaptation after it (\S\ref{sec:mitigation}).
\end{enumerate}

% ============================================================================
\section{Related Work}
\label{sec:related}

\paragraph{Corruption robustness.} ImageNet-C/CIFAR-C \cite{hendrycks2019benchmarking}
established corruption benchmarks; their brightness transform \emph{adds}
light, modeling overexposure. Low-light-specific studies typically benchmark
enhancement networks \cite{wei2018deep} rather than frozen self-supervised
representations. We bridge the two with a severity ladder matched to the
photography literature.

\paragraph{Self-supervised representation analysis.} CKA
\cite{kornblith2019similarity} is the standard tool for representational
comparison; we use its \emph{unbiased} variant, since the biased estimator
approaches 1 for independent representations at $n \ll d$ (a pitfall we
verify empirically in our released tests). Work on DINO-family features shows
they emphasize low-frequency content \cite{ning2024analyzing}, and that
training objective shapes where distribution shift lands
\cite{snell2024objectives}. Patch-embedding fragility under input shifts was
proposed in \cite{xiao2021improved}.

\paragraph{Parameter-efficient adaptation.} LoRA \cite{hu2021lora}
low-rank adapters are widely assumed to transfer robustly. Our negative
result connects to findings that adaptation objectives, not adapter placement,
determine distribution-shift behavior \cite{snell2024objectives}.

% ============================================================================
\section{Experimental Design}
\label{sec:design}

\paragraph{Corruption model.} Given a clean image $I$, severity $s \in
\{0..5\}$ defines
\begin{equation}
\label{eq:corruption}
I_{\text{obs}} = \mathrm{clip}_{[0,255]}\!\big(\mathrm{trunc}_8(a_s I + \eta)\big),
\quad \eta \sim \mathcal{N}(0,\sigma_s^2)
\end{equation}
with $(a_s)$ linearly spaced $1.0 \to 0.15$ and $(\sigma_s)$
$0 \to 13$ (Tab.~\ref{tab:ladder}). The model has two stages: the
\emph{analog} stage $a_s I + \eta$ is invertible in principle (scale by
$1/a_s$); the \emph{digitization} stage $\mathrm{trunc}_8(\cdot)$ is not.
Severity 0 is an exact identity; noise is absent, making the quantization
decomposition clean. Note Eq.~\ref{eq:corruption} \emph{truncates} (as real
ADCs and \texttt{astype(uint8)} do), rather than rounds.

\paragraph{Protocol.} CIFAR-10 \cite{krizhevsky2009learning}: 1{,}000 seeded
test images, 70/30 stratified split. DINOv2 ViT-S/14 \cite{oquab2023dinov2}
frozen; embeddings upsampled $32{\to}224$ with ImageNet normalization. A
multinomial linear probe is trained \emph{once} on clean embeddings and
evaluated unchanged at every severity --- measuring zero-shift robustness,
not adaptation. All experiments share one configuration file, hashed into the
released manifest; every table in this paper regenerates from
\texttt{run\_experiments.py}. Statistics: 2{,}000-replicate bootstrap and
Wilson intervals, exact paired permutation tests (5{,}000 replicates), Benjamini--Hochberg
correction across test families \cite{benjamini1995controlling}.

\begin{table}[t]
\centering\small
\caption{Severity ladder of the two-stage low-light corruption.}
\label{tab:ladder}
\begin{tabular}{cccl}
\toprule
$s$ & $a_s$ & $\sigma_s$ & Regime \\
\midrule
0 & 1.00 & 0  & clean reference \\
1 & 0.75 & 2  & dim evening \\
2 & 0.55 & 4  & heavy dusk \\
3 & 0.38 & 6  & deep twilight \\
4 & 0.25 & 9  & near-dark \\
5 & 0.15 & 13 & moonlight \\
\bottomrule
\end{tabular}
\end{table}

% ============================================================================
\section{The Failure Curve is Structured, Not Gradual}
\label{sec:curve}

\TableRef{main} reports accuracy and embedding drift. Three regimes emerge:
a \emph{grace} regime (sev 0--1) where accuracy drops $<$1\,pp while cosine
similarity to clean embeddings falls from 1.00 to \CosSevOne{}; a \emph{cliff}
(sev 2--4) where accuracy collapses \AccCliffDroppp{} in three steps; and a
\emph{floor} (sev 5) at \AccFloor{} --- chance level for 10 classes.

The curve is \emph{statistically non-linear}: against a linearity null, the
observed residual structure exceeds 95\% of parametric-bootstrap replicates
($p$ = \PNotNull; against midpoint symmetry, $p$ = \PMidNotNull). Accuracy is
\emph{not} a linear read of embedding drift: cosine similarity declines
smoothly while the decision collapses catastrophically. This dissociation
motivates the margin analysis of \S\ref{sec:mech}.

\begin{table}[t]
\centering\small
\caption{Main curve: linear-probe accuracy (95\% bootstrap CI), embedding
drift, and mean logit margin vs severity. Generated from
\texttt{results/main\_curve.csv}.}
\label{tab:main}
\resizebox{\linewidth}{!}{%
\begin{tabular}{cccccc}
\toprule
Sev. & $a$ & Acc (\%) & 95\% CI & cos to clean & margin \\
\midrule
0 & 1.0 & 93.3 & [88.3, 97.5] & 1.00 & 5.80 \\
1 & 0.75 & 95.0 & [90.0, 98.3] & 0.94 & 6.01 \\
2 & 0.55 & 90.0 & [84.2, 95.0] & 0.80 & 5.29 \\
3 & 0.38 & 61.7 & [52.5, 70.0] & 0.60 & 3.55 \\
4 & 0.25 & 23.3 & [15.0, 30.0] & 0.31 & 2.25 \\
5 & 0.15 & 10.8 & [5.0, 16.7] & 0.16 & 3.05 \\
\bottomrule
\end{tabular}
}

\end{table}

\begin{figure}[t]
\centering
\includegraphics[width=\linewidth]{figures/fig2_main_curve.png}
\caption{Accuracy (blue, left axis; shaded 95\% CI) and mean cosine
similarity to clean embeddings (red, right axis) vs severity. The two curves
dissociate: drift is smooth, the decision is not.}
\label{fig:main}
\end{figure}

% ============================================================================
\section{Where the Representation Breaks}
\label{sec:mech}

\paragraph{Layer-wise drift.} With unbiased linear CKA
\cite{kornblith2019similarity} between clean and degraded CLS activations per
block, drift is \emph{back-loaded}: blocks 8--11 show a mean clean$\to$darkest
CKA drop of \LateDrift{}, versus \EarlyDrift{} for blocks 0--3. The
patch-embedding-bias hypothesis \cite{xiao2021improved} predicts the opposite
gradient; the attention/norm hypothesis \cite{snell2024objectives} predicts
what we observe (Fig.~\ref{fig:cka}).

\begin{figure}[t]
\centering
\includegraphics[width=\linewidth]{figures/fig3_cka.png}
\caption{Unbiased linear CKA between clean and degraded block outputs.
Drift concentrates in late blocks (8--11).}
\label{fig:cka}
\end{figure}

\paragraph{Frequency dependence.} Pure low-pass corruption (which destroys
only high frequencies) degrades accuracy far less than pure high-pass at
matched severity, and low light --- which damages both bands jointly plus SNR
--- is significantly worse than either (paired permutation tests,
all BH-corrected $p < 0.01$). This confirms the low-frequency bias of
DINO-family features \cite{ning2024analyzing} and explains why contrast
restoration (which restores low-frequency structure) is the right lever
(\S\ref{sec:mitigation}).

\paragraph{Margins, logit lens, and token ablation.} The probe's mean logit
margin (top1$-$top2) shrinks \emph{before} accuracy collapses and is
near-zero on correctly-classified examples by sev 4 --- the cliff is a
threshold effect on the decision, not a smooth decay of feature quality. A
logit-lens over patch tokens (renormalizing the top-$k$ highest-norm token
embeddings and re-probing) shows high-norm tokens at high severity carry
\emph{misleading} evidence: ablating them (AOPC-style \cite{erichson2019aopc})
\emph{increases} accuracy relative to the full dark embedding, evidence that
darkness converts informative tokens into corrupted evidence rather than
merely deleting information.

% ============================================================================
\section{Quantization is a First-Class Cause}
\label{sec:quant}

Because Eq.~\ref{eq:corruption} decomposes into an invertible analog stage and
an irreversible digitization stage, we can measure each arm's accuracy
separately: the \emph{float arm} feeds $a_s I + \eta$ (float32) directly to
the backbone with matched preprocessing; the \emph{uint8 arm} is the standard
pipeline. \TableRef{quant} reports the decomposition. The irreversible gap
grows with severity as the signal occupies fewer of the 256 levels
(at $s{=}5$, $0.15\,I$ maps most pixels into ${\le}39$ levels, with noise
$\sigma{=}13$ scrambling those). \emph{Float-sensor inference} --- keeping
stage-1 output in float through the backbone --- recovers a substantial
fraction of the accuracy that 8-bit input loses, with zero training. To our
knowledge this is the first direct measurement of quantization's share of
low-light failure in a self-supervised backbone.

\begin{table}[t]
\centering\small
\caption{Quantization decomposition: float stage-1 vs standard uint8 pipeline.}
\label{tab:quant}
\begin{tabular}{ccccc}
\toprule
Sev. & Float (\%) & uint8 (\%) & Gap (pp) \\
\midrule
0 & 93.3 & 93.3 & +0.0 \\
1 & 94.2 & 94.2 & +0.0 \\
2 & 87.5 & 87.5 & +0.0 \\
3 & 65.0 & 63.3 & +1.7 \\
4 & 25.0 & 25.0 & +0.0 \\
5 & 10.8 & 10.8 & +0.0 \\
\bottomrule
\end{tabular}

\end{table}

% ============================================================================
\section{Adaptation Hurts; Restoration Helps}
\label{sec:lora}

The CKA analysis predicts late attention blocks are where adaptation budget
should go. We test that prediction directly: LoRA \cite{hu2021lora} on
\texttt{attn.qkv} and \texttt{attn.proj} ($r \in \{4,8,16\}$, plus MLP targets),
5{,}000 CIFAR-10 images with 70\% low-light augmentation, 10 epochs. The
result is a clean negative: \emph{every} configuration underperforms the
frozen baseline at high severity (\TableRef{lora}), and the drift-reduction
prediction fails --- adapters reduce CKA drift in some layers while reducing
accuracy.

We attribute this to objective conflict: with a 5{,}000-image, 10-way head,
the adapters receive gradients that teach CIFAR-10 class discrimination, not
photometric invariance; the low-rank update trades general features for
task-specific ones. This reframes the ``which layers to adapt'' question: no
placement choice in our grid rescued robustness, consistent with objective-dominance
findings \cite{snell2024objectives}.

\begin{table}[t]
\centering\small
\caption{LoRA rank $\times$ target ablation vs the frozen baseline. Negative
$\Delta$ = worse than doing nothing.}
\label{tab:lora}
\input{results/table_lora}
\end{table}

\label{sec:mitigation}
In contrast, three \emph{zero-training} mitigations applied to the corrupted
input --- linear gain $I/a_s$, gamma correction, and CLAHE on the LAB
L-channel \cite{zuiderveld1994adaptive} --- all improve the frozen model's
accuracy at mid severities (\TableRef{mit}), none require gradients, and all
fail at sev 5 exactly where \S\ref{sec:quant} shows information is destroyed.
The contrast between Table~\ref{tab:lora} and Table~\ref{tab:mit} is the
paper's actionable core: \emph{fix the input, don't touch the backbone}.

\begin{table}[t]
\centering\small
\caption{Zero-training mitigation vs no mitigation (frozen model).}
\label{tab:mit}
\begin{tabular}{ccccc}
\toprule
Sev. & None (\%) & Gain (\%) & Gamma (\%) & CLAHE (\%) \\
\midrule
0 & 93.3 & 93.3 & 92.5 & 65.8 \\
1 & 95.0 & 94.2 & 92.5 & 53.3 \\
2 & 90.0 & 89.2 & 86.7 & 31.7 \\
3 & 61.7 & 68.3 & 68.3 & 21.7 \\
4 & 23.3 & 30.0 & 19.2 & 11.7 \\
5 & 10.8 & 11.7 & 11.7 & 10.0 \\
\bottomrule
\end{tabular}

\end{table}

% ============================================================================
\section{Discussion: A Unifying Account}
\label{sec:discussion}

The measurements compose into one causal chain: photon loss degrades SNR and
contrast; the probe's signal --- predominantly low-frequency --- weakens;
attention over less-discriminable tokens reorganizes (late-block CKA drift);
margins shrink; the decision crosses thresholds and accuracy cliffs. Digitization
then \emph{permanently} deletes the evidence that restoration would need, placing
a hard ceiling on every mitigation. The ceiling is measurable (float arm),
the floor is verifiable (sev 5 at chance), and the middle is restorable ---
but only before quantization, not after.

Three testable corollaries: (1) accuracy tracks the margin distribution, not
mean embedding drift; (2) no mitigation can beat the float arm --- that is the
information-theoretic ceiling for this corruption; (3) adaptation budget
should target objectives (invariance losses, anchored adapters), not layer
placement.

% ============================================================================
\section{Limitations}
\label{sec:limitations}

Single backbone scale (ViT-S/14); single dataset (CIFAR-10 at 32px upsampled
to 224, which itself alters frequency content); synthetic corruption with a
single global brightness factor (no vignetting, auto-exposure, or real sensor
MTF); 300-image test fold (mitigated by exact intervals); LoRA results at one
seed per configuration (the negative effect size is consistent across all six
configurations, but per-config variance is unmeasured). Real low-light
validation (ExDark \cite{wei2018deep}, SID \cite{chen2018learning}) and
ViT-B/14 replication are released as config-flagged extensions.

% ============================================================================
\section{Conclusion}

Low-light failure in self-supervised ViTs is a readout-and-quantization
problem before it is a fine-tuning problem. The representation's low-frequency
backbone survives contrast loss but not the ADC; adaptation placed at the
drifting layers makes things worse because the objective, not the location,
is wrong; and classical restoration --- the baseline the field skips ---
recovers real accuracy for zero cost. We release the full harness, statistical
tests, and configuration manifests to make every number in this paper
reproducible from one command.

% ============================================================================
{\small
\bibliographystyle{ieee_fullname}
\bibliography{references}
}

\appendix
\section{Reproducibility}
Every experiment is a subcommand of \texttt{run\_experiments.py}
(\texttt{main\_curve}, \texttt{cka}, \texttt{frequency}, \texttt{mechanism},
\texttt{quantization}, \texttt{lora\_ablation}, \texttt{mitigation}) driven by
\texttt{configs/experiments.yaml}; the config SHA-256, package versions, and
dataset indices are written to \texttt{results/manifest.jsonl} per run. The
statistical machinery is covered by 22 unit tests (\texttt{tests/run\_tests.py}),
including verification that the biased CKA estimator inflates similarity at
$n \ll d$ while the unbiased estimator used here does not.

\end{document}
